% amsthm must be loaded before the class: sn-jnl defines its theorem
% styles (thmstyleone/two/three) inside an \@ifpackageloaded{amsthm}
% guard evaluated at class-load time.
\RequirePackage{amsthm}
\documentclass[sn-mathphys,NameDate]{sn-jnl}

%% ============================================================
%% PAPER B (JoS) — assembled 2026-06-11.
%%
%% Organizing frame (decided after the 2026-06-10 falsification):
%% "How far does permutation-schedule dominance extend beyond the
%% two-stage assembly flow shop, and where exactly does it break?"
%% Headline: the depth characterization (no directed 3-path <=> full
%% synchronization is makespan-safe for all works), which REPLACES
%% the diamond/FJFJ conjecture of the Zenodo preprint (corrected
%% here, Section "Computational falsification").
%%
%% Sources ported:
%%   [P1] paper/extended/synchronization.tex Section 5 (chains).
%%   [P2] paper/concordant/concordant-flowtime-proof.tex (full n=2
%%        proof; concordant-flowtime-general.tex is quarantined and
%%        was NOT used).
%%   [P3] paper/boundary/depth-characterization.tex (lemmas,
%%        embedding theorem, characterization, evidence).
%%   [P4] Paper A cited as bowmanA (submitted to JoS; Zenodo doi).
%%
%% Proof-hardening decisions (2026-06-11), relative to the research
%% note paper/boundary/depth-characterization.tex:
%%  (1) Prefix-queue comparison step: proved via a first-crossing
%%      argument (phi = max(S - S~, 0)). The note's statement was
%%      also CORRECTED: monotonicity/Lipschitz hold for the prefix
%%      completion maxima M_p, NOT for individual completion times
%%      (per-job monotonicity is false: delaying a high-priority
%%      job's arrival can let a low-priority job escape preemption
%%      and finish earlier). The ordered-arrival clause (max-plus
%%      recurrence) supplies the per-job statements actually needed.
%%  (2) Padding for general n: proved. Constructed schedule puts
%%      the two heavy jobs at the top of every priority order; the
%%      synchronized lower bounds touch only the heavy pair, so they
%%      survive arbitrary interleaving of padding jobs.
%%  (3) Fluid adaptation of Paper A (positive direction): proved as
%%      MAKESPAN dominance. The note's claim of pointwise domination
%%      at sinks is FALSE in the fully synchronized fluid model
%%      (counterexample in Remark rem:pointwise; Paper A's pointwise
%%      theorem inherits the downstream sequence, which full
%%      synchronization does not). The characterization only needs
%%      the makespan form, so the headline is unaffected.
%%  (4) Adversarial audit (2026-06-11): added the optimum lower
%%      bound (>= 7M - delta, mixed-priority case analysis) to
%%      Thm. thm:negative so the 1/8 limit is two-sided ("tends to"
%%      was previously only proved as liminf >= 1/8); extended
%%      Prop. chain-failure to L > 3 by epsilon-padding; chain-sweep
%%      parameters corrected to k = 2 (k = 3 rescales by 2/3, same
%%      counts -- verified by rerun: 1,414 / 1,016 per k); concordant
%%      sweep counts re-verified and restated (784 strict L=2 +
%%      21,952 strict L=3 + 46,656 with ties, zero failures; sole
%%      n=3 L=2 CE (4,5,6)/(1,9,10) confirmed); sharpness remark now
%%      uses an unweighted-flowtime discordant CE (13.5 vs 11.5).
%%  (5) Integration of verified computational results (2026-06-11,
%%      commit c35d980 evidence): practice section gains the
%%      chain3/FJFJ policy-benchmark families (Table tab:policy)
%%      and the 561,633-instance per-pi policy==full-sync identity
%%      (incl. two-source FJFJ); robustness extended to chain3
%%      (3.17%) and FJFJ (4.11%); Discussion's integer-model
%%      "in progress" flag replaced by the settled finding
%%      (failures survive: w=[(1,3),(4,1),(1,3)] k=2 sync 6 vs
%%      opt 5, 2/4,096; knife-edge: gap 0 at scales 2,3,4).
%%      Two-sink pointwise 35.7% qualified as per-sink. Every
%%      number re-verified against an independent reimplementation
%%      before inclusion.
%% ============================================================

\usepackage{graphicx}
\usepackage{multirow}
\usepackage{amsmath,amssymb,amsfonts}
\usepackage{amsthm}
\usepackage{mathrsfs}
\usepackage{xcolor}
\usepackage{textcomp}
\usepackage{manyfoot}
\usepackage{booktabs}

\theoremstyle{thmstyleone}
\newtheorem{theorem}{Theorem}
\newtheorem{lemma}[theorem]{Lemma}
\newtheorem{corollary}[theorem]{Corollary}
\newtheorem{proposition}[theorem]{Proposition}
\theoremstyle{thmstylethree}
\newtheorem{definition}[theorem]{Definition}
\theoremstyle{thmstyletwo}
\newtheorem{remark}[theorem]{Remark}

\raggedbottom

\begin{document}

\title[The Boundary of Synchronization Dominance]{The Boundary of
Synchronization Dominance: A Depth Characterization for
Divisible-Work Scheduling on DAGs}

\author*[1]{\fnm{Eric} \sur{Bowman}}\email{eric.bowman@king.com}

\affil[1]{\orgname{King}, \orgaddress{\city{Berlin}, \country{Germany}}}

\abstract{Permutation schedules pointwise-dominate in the
divisible-work two-stage assembly flow shop. We determine
exactly how far this extends when the two-stage structure
generalizes to chains and to arbitrary DAGs of stages under a
preemptive divisible-work model. Our main result is a complete
characterization: fully synchronized schedules---one common
priority order at every stage---are makespan-optimal at every
sink for every assignment of work if and only if the DAG
contains no directed path through three stages. The positive
direction reduces, sink by sink, to a fluid form of the
two-stage capacity argument; the negative direction embeds a
three-stage chain counterexample along any longer path via a
scaling argument, with relative optimality gap approaching
$1/8$ of the synchronized makespan. The characterization
corrects a conjecture from our earlier preprint: makespan
dominance was conjectured for diamond and
fork-join--fork-join topologies on computational evidence
whose sweeps held source works fixed; with source works
adversarial, failures are in fact common. For chains we show
makespan dominance holds at length two and fails from length
three; for total flowtime we prove dominance for two jobs with
concordant work profiles and show it fails beyond. Finally we
quantify the practical stakes: a composite policy that
synchronizes only the source stages and serves downstream work
in arrival order provably reproduces the synchronized schedule
on single-source DAGs, is within $0.06\%$ of the enumerated
optimum on average over an exhaustive instance family, retains
most of its value under $\pm 50\%$ work-estimation error, and
the residual gap concentrates on rare, identifiable
cross-stage work inversions.}

\keywords{assembly flow shop, permutation schedule, divisible
work, DAG scheduling, synchronization, makespan}

\pacs[MSC Classification]{90B35}

\maketitle

%% ============================================================
\section{Introduction}
\label{sec:intro}

For the two-stage assembly flow shop with divisible work---$m$
parallel machines of $k$ identical workers feeding a single
assembly stage---permutation schedules dominate pointwise: for
every schedule there is a schedule in which all stage-1
machines share one job order and every job's assembly
completion is no later \citep{bowmanA}. The result is
constructive, holds for every regular objective, and is
formalized in Lean~4. It invites an obvious question: how much
of the organization-like structure around the two-stage core
can the dominance survive? Deeper pipelines? Intermediate
parallel layers? Branching without reconvergence?

This paper answers the question completely for the makespan
criterion in the preemptive divisible-work model. The answer
is a sharp structural boundary:

\begin{quote}
\emph{Full synchronization---one common priority order at
every stage---is makespan-safe for every work assignment if
and only if no directed path of the stage DAG passes through
three stages.}
\end{quote}

The ``if'' direction is a fluid form of the two-stage capacity
argument applied sink by sink
(Theorem~\ref{thm:positive}). The ``only if'' direction is
new: any three-stage path admits an embedding of a chain
counterexample, made robust to the surrounding DAG by a
scaling argument, and the optimality gap approaches $1/8$ of
the synchronized makespan (Theorem~\ref{thm:negative}). The
mechanism is \emph{overtaking}: when a job is light at one
stage and heavy at the next, optimal schedules let it pass its
neighbour between stages; a common order forbids this at every
depth.

The characterization corrects the record of our earlier
preprint \citep{bowmanA}, which conjectured---from more than
$2 \times 10^7$ verified instances---that makespan dominance
survives on diamond and fork-join--fork-join topologies,
attributing the apparent rescue to intermediate parallel
width. The conjecture is false, and the evidence was an
artifact of sweep design: the exhaustive sweeps varied work
only at non-source stages. Embedding the chain counterexample
with adversarial \emph{source} works produces failures in all
of these topologies, at rates near $5\%$ of random instances;
an exhaustive small-instance sweep including source works
finds failures at $0.55\%$ density
(Section~\ref{sec:falsification}). Intermediate width neither
helps nor hurts: depth alone decides.

Around the makespan boundary we assemble the adjacent
dominance landscape. For chains of stages, makespan dominance
holds at length two and fails from length three
(Section~\ref{sec:chains}). For total flowtime on chains,
synchronization is optimal for two jobs whose work profiles
are \emph{concordant}---the lighter job is lighter at every
stage---and fails beyond, both without concordance and with
three concordant jobs (Section~\ref{sec:flowtime}).

Finally, because the negative results say only that full
synchronization is not \emph{free} at depth three, we quantify
what it actually costs and what a manager should do instead
(Section~\ref{sec:practice}). A composite policy---one shared
priority order at the source stages, first-in-first-out by
arrival everywhere downstream---provably reproduces the fully
synchronized schedule on single-source DAGs
(Lemma~\ref{lem:propagation}); over an exhaustive family of
$65{,}536$ instances it is within $0.06\%$ of the enumerated
optimum on average and optimal outright in $99.4\%$ of
instances, with the same pattern across chain and two-source
fork-join families; and its near-optimality survives
$\pm 50\%$ work-estimation error. Where full synchronization
does fail, the gap is material (mean $5.5$--$10.4\%$ across
families) but rare ($0.55$--$10\%$ of instances) and
identifiable: it requires a cross-stage work inversion, which
deliberate expediting can capture.

%% ============================================================
\section{Related work}
\label{sec:related}

The two-stage assembly flow shop was introduced by
\citet{lee1993}. \citet{potts1995} proved strong NP-hardness
for $m \geq 2$ and the foundational dominance result:
permutation schedules suffice for makespan, reducing the
search space from $(n!)^m$ to $n!$. \citet{tozkapan2003}
extended permutation dominance to total weighted flowtime;
\citet{hariri1997} and \citet{hadda2024} developed
branch-and-bound methods built on dominance rules.
\citet{koulamas2007} studied uniform parallel machines at
stage~1. Our companion paper \citep{bowmanA} extends the
dominance pointwise to a divisible-work model with $k$
identical workers per machine and machine-dependent works;
the present paper takes that result as its starting point and
maps its outer boundary.

Adjacent literatures do not address the question. The
malleable- and divisible-load literature
\citep{drozdowski2009} studies divisible work on parallel
processors without the assembly or DAG-precedence structure
between stages. Distributed assembly permutation flow shops
\citep{hatami2013} concern multiple factories rather than
per-stage priority synchronization. \citet{zhang2022} schedule
assembly with parallel worker teams in a multi-objective
setting without dominance results. The survey of
\citet{komaki2019} covers assembly flow shop variants; neither
the divisible-work model nor the question of when a single
shared priority order is lossless across a DAG of stages
appears there. \citet{alanzi2007} use permutation dominance to
restrict heuristic search in the two-stage setting with
setups. In the three-field notation of \citet{graham1979} our
positive results live in assembly-type extensions of
$F2 \mid \mathrm{div} \mid C_{\max}$; the negative results
concern arbitrary stage DAGs under a preemptive fluid relaxation.

%% ============================================================
\section{The fluid model and basic lemmas}
\label{sec:model}

\begin{definition}[Fluid preemptive priority DAG model]
\label{def:model}
A DAG $G = (V, E)$ of \emph{stages}; stage $v$ has $k_v \geq 1$
identical workers and work $w_{v,j} > 0$ for each job
$j \in [n]$. Job $j$ \emph{arrives} at $v$ at
$a_v(j) = \max_{u : (u,v) \in E} C_u(j)$ (at sources,
$a_v(j) = 0$). A schedule assigns each stage a priority order
$\beta_v \in \mathfrak{S}_n$; at every instant, stage $v$
devotes all $k_v$ workers (total rate $k_v$) to the
highest-$\beta_v$-priority job that has arrived and is
incomplete at $v$. $C_v(j)$ denotes the completion time of $j$
at $v$. A schedule is \emph{synchronized} (a permutation
schedule) if $\beta_v = \pi$ for all $v$. For a sink $t$, the
makespan at $t$ is $\max_j C_t(j)$. We say \emph{makespan
dominance holds at sink $t$} for an instance if
\[
  \min_{\pi \in \mathfrak{S}_n} \max_j C^{\pi}_t(j)
  \;=\;
  \min_{\boldsymbol{\beta}} \max_j C^{\boldsymbol\beta}_t(j).
\]
\end{definition}

Each stage's trajectory is piecewise linear with breakpoints
only at arrivals and completions, so all completion times are
well defined and finite. Two elementary properties of the
dynamics: they are \emph{time-invariant} (shifting all
arrivals at a stage by $\varepsilon$ shifts its completions by
exactly $\varepsilon$) and \emph{positively homogeneous}
(scaling all works by $M$ scales all completion times
by~$M$). Since a stage's dynamics depend on works only through
the normalized quantities $w_{v,j}/k_v$, we refer to these as
\emph{processing times} and state instances either way.

The workhorse is a reduction of each priority class to a
single-server fluid queue.

\begin{lemma}[Prefix-queue reduction]
\label{lem:prefix}
Fix a stage with priority order $\beta$, rate $k$, works
$w_j > 0$, and arrival times $a(j)$. For $p \in [n]$ let
$P_p = \{\beta(1), \ldots, \beta(p)\}$,
$W_p = \sum_{j \in P_p} w_j$, and let
$A_p(t) = \sum_{j \in P_p,\, a(j) \leq t} w_j$ be the arrived
prefix work. Let $S_p(t)$ denote the cumulative service
delivered to jobs of $P_p$ by time $t$. Then:
\begin{enumerate}
\item[\textup{(i)}] $S_p$ is continuous, $S_p(0) = 0$,
  $S_p(t) \leq A_p(t)$ for all $t$, and between breakpoints
  \[
    \tfrac{d}{dt} S_p(t) \;=\; k \cdot
    \mathbf{1}\bigl[A_p(t) > S_p(t)\bigr].
  \]
\item[\textup{(ii)}] The time at which all of $P_p$ is
  complete satisfies
  \[
    M_p \;:=\; \max_{j \in P_p} C(j)
    \;=\; \inf\{\, t : S_p(t) = W_p \,\}.
  \]
\item[\textup{(iii)}] \textup{(Invariance.)} The stage
  makespan $M_n$ depends only on the arrival times and works,
  not on $\beta$.
\item[\textup{(iv)}] \textup{(Monotonicity and Lipschitz
  bound.)} Each $M_p$ is non-decreasing in every arrival time,
  and $|M_p(a') - M_p(a)| \leq \sup_j |a'(j) - a(j)|$.
\item[\textup{(v)}] \textup{(Ordered arrivals.)} If
  $a(\beta(1)) \leq \cdots \leq a(\beta(n))$, then no job is
  ever preempted, $C(\beta(p)) = M_p$ for every $p$, and the
  completions obey the max-plus recurrence
  \[
    C(\beta(p)) = \max\bigl(a(\beta(p)),\,
      C(\beta(p-1))\bigr) + w_{\beta(p)}/k,
    \qquad C(\beta(0)) := 0 .
  \]
\end{enumerate}
\end{lemma}

\begin{proof}
(i) Whenever some job of $P_p$ has arrived and is incomplete,
the highest-priority arrived incomplete job lies in $P_p$, so
the stage serves a job of $P_p$ at full rate $k$; otherwise it
serves outside $P_p$ or idles, and $S_p$ is flat. The
condition ``some job of $P_p$ arrived and incomplete'' is
exactly $A_p(t) > S_p(t)$, and service cannot exceed arrived
work, giving $S_p \leq A_p$.

(ii) All of $P_p$ is complete at time $t$ if and only if all
arrived prefix work has been served and no prefix work remains
to arrive, i.e.\ $S_p(t) = W_p$; the infimum is attained by
continuity of $S_p$.

(iii) For $p = n$, $A_n$ is the total arrived work, which does
not depend on $\beta$, and by (i) the trajectory $S_n$ is
determined by $A_n$; by (ii), $M_n$ is determined by $S_n$.

(iv) Let $a \leq a'$ componentwise, with arrived-work curves
$A_p \geq A'_p$ pointwise, and trajectories $S_p$, $S'_p$ from
(i). Set $\varphi(t) = \max\bigl(S'_p(t) - S_p(t),\, 0\bigr)$,
a continuous piecewise-linear function with $\varphi(0) = 0$.
At any non-breakpoint $t$ with $S'_p(t) > S_p(t)$ we have
$A_p(t) \geq A'_p(t) \geq S'_p(t) > S_p(t)$, so
$\frac{d}{dt}S_p(t) = k \geq \frac{d}{dt}S'_p(t)$ and
$\varphi'(t) \leq 0$; where $S'_p \leq S_p$, $\varphi$ is
locally zero. Hence $\varphi \equiv 0$, i.e.\
$S'_p \leq S_p$ pointwise, and by (ii)
$M_p(a) \leq M_p(a')$. The Lipschitz bound follows from
monotonicity and time-invariance: if
$a \leq a' \leq a + \varepsilon\mathbf{1}$ then
$M_p(a) \leq M_p(a') \leq M_p(a) + \varepsilon$, and the
general case reduces to this with
$\varepsilon = \sup_j |a'(j) - a(j)|$ after replacing $a$ by
the componentwise minimum.

(v) With $\beta$-ordered arrivals, whenever $\beta(p)$ is
present, every incomplete higher-priority job has already
arrived, so $\beta(p)$ is served only when all of $P_{p-1}$ is
complete, and from that moment it is served at rate $k$
without interruption. Hence $\beta(p)$ starts at
$\max\bigl(a(\beta(p)), C(\beta(p-1))\bigr)$
(using $C(\beta(p-1)) = M_{p-1}$ inductively), runs for
$w_{\beta(p)}/k$, and completes last among $P_p$, giving
$C(\beta(p)) = M_p$ and the recurrence.
\end{proof}

Part (iv) is deliberately stated for the prefix maxima $M_p$
and not for individual completion times: individual
completions are \emph{not} monotone in arrivals (delaying a
high-priority job's arrival can let a low-priority job escape
preemption and finish earlier). Part (v) recovers per-job
control exactly when arrivals respect the priority order,
which is the regime our positive arguments arrange.

\begin{lemma}[Light stages perturb little]
\label{lem:light}
Let stage $v$ have total work
$\overline{W}_v = \sum_j w_{v,j}$. Then for every priority
order and every job $j$:
\textup{(i)} $C_v(j) \leq a_v(j) + \overline{W}_v / k_v$, and
\textup{(ii)} $\max_j C_v(j) \leq \max_i a_v(i) +
\overline{W}_v / k_v$.
\end{lemma}

\begin{proof}
(i) From $a_v(j)$ until $C_v(j)$, job $j$ is present and
incomplete, so the stage continuously serves \emph{some} job
at rate $k_v$; total work is $\overline{W}_v$, so the busy
time is at most $\overline{W}_v/k_v$. (ii) After
$\max_i a_v(i)$ all work has arrived; work conservation
exhausts it within $\overline{W}_v/k_v$.
\end{proof}

%% ============================================================
\section{Chains: dominance to depth two, failure from depth
three}
\label{sec:chains}

A \emph{chain} of length $L$ is the path DAG
$1 \to 2 \to \cdots \to L$ in
Definition~\ref{def:model}: stage $\ell$ has $k_\ell \geq 1$
workers and works $w_{\ell,j}$, all jobs are released to
stage~1 at time~0, and stage $\ell \geq 2$ can begin job $j$
only after stage $\ell - 1$ completes it. This is the
two-stage assembly model with $m = 1$ when $L = 2$, and the
simplest setting in which depth can be isolated from width.

\begin{lemma}[Priority-order completion]
\label{lem:priority-order}
Under a synchronized schedule with order $\pi$ on a chain,
jobs complete each stage in priority order:
$C_\ell(\pi(1)) \leq \cdots \leq C_\ell(\pi(n))$ for every
stage $\ell$.
\end{lemma}

\begin{proof}
By induction on $\ell$. At stage~1 all jobs arrive at time~0,
which is trivially $\pi$-ordered; at stage $\ell \geq 2$ the
arrivals are the stage-$(\ell-1)$ completions, $\pi$-ordered
by the inductive hypothesis. In both cases
Lemma~\ref{lem:prefix}(v) applies and yields $\pi$-ordered
completions.
\end{proof}

\begin{theorem}[Two-stage chain makespan dominance]
\label{thm:chain-l2}
For $L = 2$, the optimal makespan over all per-stage orderings
equals the optimal makespan over synchronized schedules.
\end{theorem}

\begin{proof}
Consider any schedule with orderings $\beta_1, \beta_2$. By
Lemma~\ref{lem:prefix}(iii), the makespan at stage~2 depends
only on the arrival times at stage~2 (the stage-1 completions)
and the stage-2 works, not on $\beta_2$. Therefore the
synchronized schedule $(\beta_1, \beta_1)$ achieves the same
makespan as $(\beta_1, \beta_2)$. Since this holds for every
$\beta_1$, the minimum over synchronized schedules equals the
minimum over all schedules.
\end{proof}

For longer chains the natural inductive extension---apply the
$(L-1)$-stage result, then stage-$L$ invariance---requires
that an ordering good for stages $1, \ldots, L{-}1$ is also
good for stage~$L$. This fails when work profiles invert
across stages, and the failure is not an artifact of exotic
objectives: it already defeats the makespan.

\begin{proposition}[Failure of chain dominance for $L \geq 3$]
\label{prop:chain-failure}
For $L \geq 3$, there exist instances where no synchronized
schedule achieves the optimal makespan, and instances where no
synchronized schedule achieves the optimal total weighted
completion time.
\end{proposition}

\begin{proof}
Counterexample: $L = 3$, $n = 2$, $k_\ell = 2$,
$w = [(1,2),\,(4,1),\,(1,2)]$, i.e.\ processing times
$(0.5, 1.0)$, $(2.0, 0.5)$, $(0.5, 1.0)$ for jobs $(1,2)$
stage by stage.

Under the synchronized order $(1,2)$: at stage~1, job~1
completes at $0.5$ and job~2 at $1.5$. At stage~2, job~1
occupies $[0.5, 2.5]$; job~2 then occupies $[2.5, 3.0]$. At
stage~3, job~1 occupies $[2.5, 3.0]$ and job~2 $[3.0, 4.0]$.
Completions $(3.0, 4.0)$: makespan $4.0$, flowtime $7.0$.

Under the synchronized order $(2,1)$: stage~1 gives
$(1.5, 1.0)$; stage~2: job~2 occupies $[1.0, 1.5]$, job~1
$[1.5, 3.5]$; stage~3: job~2 $[1.5, 2.5]$, job~1
$[3.5, 4.0]$. Completions $(4.0, 2.5)$: makespan $4.0$,
flowtime $6.5$.

Under the unsynchronized orders
$(1,2) \mid (2,1) \mid \text{any}$ (the stage-3 ordering is
irrelevant for the makespan by
Lemma~\ref{lem:prefix}(iii)): stage~1 gives $(0.5, 1.5)$. At
stage~2, job~1 starts at $0.5$; job~2 arrives at $1.5$ with
higher priority and preempts it (job~1 has $1.0$ of $2.0$
processing done); job~2 completes at $2.0$; job~1 resumes and
completes at $3.0$. At stage~3 (order $(2,1)$), job~2 arrives
at $2.0$ and completes at $3.0$; job~1 arrives at $3.0$ and
completes at $3.5$. Completions $(3.5, 3.0)$: makespan $3.5$,
flowtime $6.5$.

Hence the unsynchronized schedule achieves strictly better
makespan ($3.5 < 4.0$) and equal-or-better flowtime;
enumeration over all $2^3 = 8$ per-stage orderings confirms no
schedule synchronized across the three stages achieves
makespan below $4.0$. For weighted completion time, the
stage-3 completion vectors are $C^{(1,2)} = (3.0, 4.0)$,
$C^{(2,1)} = (4.0, 2.5)$, and
$C^{\mathrm{unsync}} = (3.5, 3.0)$; with weights $v = (3,2)$
the synchronized schedules both give $17.0$ while the
unsynchronized schedule gives $16.5$.

For $L > 3$, append stages at which both jobs have work
$\varepsilon > 0$. Arrivals at each appended stage are the
previous stage's completions and completions only increase
along the chain, so under every synchronized order the final
makespan is still at least $4.0$ and the weighted completion
time at least $17.0$; extending the unsynchronized schedule
through the appended stages adds at most $2\varepsilon/k_\ell$
per stage to each completion
(Lemma~\ref{lem:light}(i)). For $\varepsilon$ small, both
strict separations persist.
\end{proof}

\begin{remark}[Mechanism: cross-stage work inversion]
\label{rem:mechanism}
The counterexample exploits preemption at an intermediate
stage. Job~1 is light at stage~1 but heavy at stage~2; job~2
is the reverse. The unsynchronized schedule uses opposite
priorities at stages~1 and~2: stage~1 completes the light
job~1 early so it flows downstream; at stage~2, job~2
overtakes it, completing its short stage-2 work immediately
and feeding stage~3 sooner. A synchronized schedule cannot
express this overtaking. Systematic search over $46{,}656$
chain configurations ($L = 3$, $n = 2$, $k = 2$, work
values $1$--$6$) found $1{,}414$ makespan counterexamples and
$1{,}016$ flowtime counterexamples; the largest makespan gap
in that range is $1.0$ (e.g.\ synchronized $5.5$ vs.\
unsynchronized $4.5$). By positive homogeneity the $k = 3$
sweep yields the same counts, with gaps scaled by $2/3$.
\end{remark}

\begin{remark}[Pointwise domination fails already at $L = 2$]
\label{rem:pointwise}
Theorem~\ref{thm:chain-l2} is about the makespan, and that is
not an accident of presentation. In the fully synchronized
sense of Definition~\ref{def:model}---the sink must use the
common order too---\emph{pointwise} domination fails already
on two-stage chains. Take $L = 2$, $n = 2$, $k_\ell = 1$,
$w_1 = (1, 2)$, $w_2 = (3, 1)$, and the schedule
$\beta_1 = (1,2)$, $\beta_2 = (2,1)$: stage-1 completions are
$(1, 3)$; stage~2 serves job~1 on $[1,3)$, is preempted by
job~2 on $[3,4]$, and finishes job~1 at $5$, giving
completions $(5, 4)$. The synchronized order $(1,2)$ gives
$(4, 5)$ (worse for job~2); the order $(2,1)$ gives $(6, 3)$
(worse for job~1). Neither dominates pointwise. This does not
contradict the pointwise theorem of \citet{bowmanA}: there the
constructed permutation schedule synchronizes the
\emph{upstream} machines and \emph{inherits} the downstream
sequence of the original schedule, whereas full
synchronization forces the downstream order as well. Makespan
dominance is the strongest form that survives full
synchronization, and it is the form our characterization
settles.
\end{remark}

%% ============================================================
\section{The boundary: a depth characterization}
\label{sec:boundary}

We now place the chain counterexample inside an arbitrary DAG.
Throughout this section, fix a directed path
$u \to v \to x$ in $G$ (three distinct stages), let
$n \geq 2$, and call jobs $1$ and $2$ the \emph{heavy} jobs
and jobs $3, \ldots, n$ \emph{padding} jobs. Given a scale
$M \in \mathbb{Z}_{>0}$, define the instance $I_M$:
\[
  w_u = (k_u M,\, 2 k_u M),\quad
  w_v = (4 k_v M,\, k_v M),\quad
  w_x = (k_x M,\, 2 k_x M)
\]
for the heavy jobs (so the path processing times are
$(M, 2M)$, $(4M, M)$, $(M, 2M)$: the instance of
Proposition~\ref{prop:chain-failure} scaled by $2M$), work $1$
for every padding job at $u, v, x$, and work
$1$ for every job at every off-path stage. Let
\[
  \delta \;:=\; n\,|V|
  \;\geq\; \sum_{y \notin \{u,v,x\}} \overline{W}_y / k_y
  \;+\; \max_{s \in \{u,v,x\}} \overline{W}^{\mathrm{pad}}_s / k_s ,
\]
a crude aggregate bound on all off-path and padding work
(normalized); here
$\overline{W}^{\mathrm{pad}}_s = n - 2$ is the padding work at
a path stage. Two consequences of Lemma~\ref{lem:light} used
repeatedly: (a) every off-path stage all of whose ancestors are
off-path completes all jobs by $\delta$ (sum
Lemma~\ref{lem:light}(ii) along ancestors); in particular all
arrivals at $u$ are at most $\delta$; (b) for any schedule and
any job $j$, a route through off-path stages delays $j$ by at
most $\delta$ in aggregate beyond its completion at the last
path stage on the route (sum Lemma~\ref{lem:light}(i)).

\begin{theorem}[Any three-stage path defeats full
synchronization]
\label{thm:negative}
Let $G$ contain a directed path $u \to v \to x$ and let the
worker counts $k_\cdot \geq 1$ be arbitrary. For every
$n \geq 2$ and every integer $M > 6\delta$, the instance
$I_M$ satisfies: at every sink $t$ reachable from $x$,
\[
  \min_{\pi} \max_j C^\pi_t(j) \;\geq\; 8M - \delta,
  \qquad
  \min_{\boldsymbol\beta} \max_j C^{\boldsymbol\beta}_t(j)
  \;\leq\; 7M + 5\delta .
\]
In particular makespan dominance fails at $t$, every
synchronized schedule exceeds the optimum by at least
$M - 6\delta$, and the relative gap
$1 - C^{\mathrm{opt}}_t / C^{\mathrm{sync}}_t$ tends to $1/8$
as $M \to \infty$.
\end{theorem}

\begin{proof}
The heavy jobs' processing times at the path stages are
$(M, 2M)$, $(4M, M)$, $(M, 2M)$: the instance of
Proposition~\ref{prop:chain-failure} scaled by $2M$, whose
isolated synchronized and optimal makespans are $8M$ and
$7M$. We use two facts valid for every priority schedule:
for all $v, j$,
\begin{itemize}
\item[(F1)] $C_v(j) \geq a_v(j) + w_{v,j}/k_v$, and
\item[(F2)] if $j'$ has lower $\beta_v$-priority than $j$,
  then $j'$ receives no service at $v$ during any interval in
  which $j$ is present and incomplete.
\end{itemize}

\emph{Synchronized lower bound.} Fix any $\pi$ and suppose
first that $\pi$ ranks job~1 above job~2. At $u$: by (F1),
$C_u(1) \geq M$. By (F2), job~2 is served at $u$ only
before $a_u(1) \leq \delta$ or after $C_u(1)$, so
$C_u(2) \geq C_u(1) + (2M - \delta)$. At $v$: arrivals
satisfy $a_v(1) \geq C_u(1)$ and, by consequence~(b) above
applied to job~1's off-path predecessors of $v$,
$a_v(1) \leq C_u(1) + \delta$; also
$a_v(2) \geq C_u(2) \geq C_u(1) + 2M - \delta
 \geq a_v(1) + 2M - 2\delta > a_v(1)$.
By (F1), $C_v(1) \geq a_v(1) + 4M \geq C_u(1) + 4M \geq 5M$.
Since job~1 is present at $v$ from $a_v(1) \leq a_v(2)$ and
incomplete until $C_v(1)$, (F2) gives job~2 no service at $v$
before $C_v(1)$, so $C_v(2) \geq C_v(1) + M \geq 6M$.
At $x$: by (F1),
$C_x(2) \geq a_x(2) + 2M \geq C_v(2) + 2M \geq 8M$.

If instead $\pi$ ranks job~2 above job~1, symmetrically:
$C_u(2) \geq 2M$; job~1 is served at $u$ only before
$a_u(2) \leq \delta$ or after $C_u(2)$, so
$C_u(1) \geq C_u(2) + M - \delta$. At $v$:
$a_v(2) \leq C_u(2) + \delta$ and
$a_v(1) \geq C_u(1) \geq a_v(2) + M - 2\delta > a_v(2)$;
by (F1), $C_v(2) \geq a_v(2) + M \geq C_u(2) + M$; by
(F2), job~1 is served at $v$ only after $C_v(2)$ (it arrives
after $a_v(2)$, and job~2 is present and incomplete until
$C_v(2)$), so
$C_v(1) \geq \max(a_v(1), C_v(2)) + 4M
 \geq C_u(2) + M - \delta + 4M \geq 7M - \delta$.
At $x$: $C_x(1) \geq a_x(1) + M \geq C_v(1) + M
\geq 8M - \delta$.

In both cases some job completes at $x$ no earlier than
$8M - \delta$. For any sink $t$ reachable from $x$, arrivals
propagate: along any directed route from $x$ to $t$, each
stage's makespan is at least its largest arrival, hence at
least $\max_j C_x(j) \geq 8M - \delta$. Note that padding jobs
appear in these arguments only through the bounds
$a_u(\cdot) \leq \delta$ and consequence~(b); their positions
in $\pi$ are immaterial, since (F1) and (F2) were applied only
to the heavy pair.

\emph{Constructed upper bound.} Use the orders
$(1, 2, \text{padding})$ at $u$, $(2, 1, \text{padding})$ at
$v$ and $x$ (padding jobs in a fixed arbitrary order below
both heavy jobs at all three stages), and any fixed order at
off-path stages. At $u$: job~1 is served on arrival, so
$C_u(1) = a_u(1) + M \leq M + \delta$; the heavy pair is
served continuously from
$\max(a_u(1), a_u(2)) \leq \delta$ until both complete, so
$C_u(2) \leq 3M + \delta$. By consequence~(b), heavy
arrivals at $v$ satisfy $a_v(j) \leq C_u(j) + \delta$. At $v$
(order $(2,1)$): $C_v(2) \leq a_v(2) + M
\leq 4M + 2\delta$; job~1 is served at $v$ from $a_v(1)$
onward except while job~2 is present and incomplete, an
interruption totalling at most $M$, so
$C_v(1) \leq a_v(1) + 4M + M \leq 6M + 2\delta$. At $x$
(order $(2,1)$): $a_x(2) \leq C_v(2) + \delta$, so
$C_x(2) \leq a_x(2) + 2M \leq 6M + 3\delta$; and
$a_x(1) \leq C_v(1) + \delta \leq 6M + 3\delta$, after which
job~1 waits at most for job~2 and then runs, so
$C_x(1) \leq \max(a_x(1), C_x(2)) + M \leq 7M + 3\delta$.
Padding jobs: at each path stage they are served once the
heavy jobs are out of the way, so all work at $u$, $v$, $x$ is
exhausted by $C_u(2) + \delta$, $C_v(1) + \delta$,
$C_x(1) + \delta$ respectively (Lemma~\ref{lem:light}(ii)
started from the later of the last arrival and the heavy
completion); hence the makespan at $x$ over all jobs is at
most $7M + 4\delta$. Downstream off-path stages to any sink
add at most $\delta$ in aggregate (consequence~(b)), so the
makespan at $t$ is at most $7M + 5\delta$.

\emph{Optimum lower bound.} The synchronized lower bound used
the priority between the heavy jobs only at $u$ and at $v$, so
it applies verbatim to any schedule whose orders at $u$ and $v$
agree on the heavy pair, giving makespan at least
$8M - \delta$ at $x$. It remains to bound the two mixed cases.
If $u$ ranks job~2 first and $v$ ranks job~1 first, then as
before $C_u(1) \geq C_u(2) + M - \delta \geq 3M - \delta$, so
$C_v(1) \geq a_v(1) + 4M \geq 7M - \delta$ and
$C_x(1) \geq C_v(1) + M \geq 8M - \delta$. If $u$ ranks job~1
first and $v$ ranks job~2 first, then
$a_v(2) \geq C_u(2) \geq C_u(1) + 2M - \delta
\geq a_v(1) + 2M - 2\delta > a_v(1)$. Job~2 is present and
incomplete at $v$ throughout $[a_v(2),\, C_v(2)]$, an interval
of length at least $M$, during which (F2) denies job~1 all
service at $v$. If $C_v(1) \leq a_v(2)$ then
$C_v(2) \geq a_v(2) + M \geq C_v(1) + M
\geq a_v(1) + 5M \geq 6M$ and
$C_x(2) \geq C_v(2) + 2M \geq 8M$; otherwise job~1 is
incomplete throughout that interval, which therefore lies
inside $[a_v(1),\, C_v(1)]$, so
$C_v(1) \geq a_v(1) + 4M + M \geq 6M$ and
$C_x(1) \geq C_v(1) + M \geq 7M$. In every case the makespan
at $x$---hence at every sink $t$ reachable from $x$---is at
least $7M - \delta$.

\emph{Conclusion.} For $M > 6\delta$ the two displayed bounds
are incompatible with dominance:
$8M - \delta > 7M + 5\delta$. For the limit, a computation
identical to the constructed upper bound but with the
synchronized order $(1, 2, \text{padding})$ everywhere shows
the best synchronized makespan at $t$ is also at most
$8M + 5\delta$, so
\[
  1 - \frac{7M + 5\delta}{8M - \delta}
  \;\leq\;
  1 - \frac{C^{\mathrm{opt}}_t}{C^{\mathrm{sync}}_t}
  \;\leq\;
  1 - \frac{7M - \delta}{8M + 5\delta},
\]
and both ends tend to $1/8$ as $M \to \infty$. \qedhere
\end{proof}

\begin{theorem}[Depth two is safe]
\label{thm:positive}
If every directed path of $G$ visits at most two stages
(equivalently, every stage is a source or a sink), then for
every work assignment, every worker assignment, and every sink
$t$, makespan dominance holds at $t$.
\end{theorem}

\begin{proof}
Fix a sink $t$ and any schedule $\boldsymbol\beta$. Stages
that do not feed $t$ cannot influence completions at $t$
(stages share no workers), so restrict attention to $t$ and
its predecessors, all of which are sources (a predecessor with
a predecessor would create a three-stage path): a two-stage
assembly instance with all jobs available at the sources at
time zero.

Let $r_j = \max_{i \in \mathrm{pred}(t)} C_i(j)$ be the
release times of $\boldsymbol\beta$, and let $\pi$ sort jobs
by non-decreasing $r_j$. \emph{Capacity bound:} in
$\boldsymbol\beta$, source $i$ completes each of
$\pi(1), \ldots, \pi(s)$ by $r_{\pi(s)}$, and its service rate
is at most $k_i$, so
$W_i(s) := \sum_{q \leq s} w_{i,\pi(q)} \leq k_i\,
r_{\pi(s)}$. Under the synchronized order $\pi$, source $i$
serves the prefix continuously from time zero
(Lemma~\ref{lem:prefix}(v) with all arrivals zero), so
$C^\pi_i(\pi(s)) = W_i(s)/k_i \leq r_{\pi(s)}$, whence the
synchronized release times satisfy $r^\pi_j \leq r_j$ for
every $j$. \emph{Sink:} by Lemma~\ref{lem:prefix}(iii) the
makespan at $t$ is a function of its arrival vector only---in
particular the original schedule's makespan at $t$ is
$M_t(r)$ regardless of $\beta_t$---and by
Lemma~\ref{lem:prefix}(iv) this function is componentwise
non-decreasing. Hence the synchronized makespan is
$M_t(r^\pi) \leq M_t(r)$. Since $\boldsymbol\beta$ was
arbitrary, the minimum over synchronized schedules equals the
minimum over all schedules.
\end{proof}

\begin{corollary}[Characterization]
\label{cor:characterization}
In the model of Definition~\ref{def:model}, the following are
equivalent for a DAG $G$:
\begin{enumerate}
\item[\textup{(i)}] for every assignment of works and workers,
  makespan dominance holds at every sink;
\item[\textup{(ii)}] $G$ contains no directed path through
  three stages.
\end{enumerate}
\end{corollary}

\begin{proof}
(ii) $\Rightarrow$ (i) is Theorem~\ref{thm:positive}. If (ii)
fails, $G$ contains a path $u \to v \to x$, and
Theorem~\ref{thm:negative} produces an instance and a sink at
which dominance fails, contradicting (i).
\end{proof}

\begin{remark}[Reading for organizations]
\label{rem:org}
Condition (ii) is precisely the ``parallel build $\to$
integrate'' shape. The moment any stage's output feeds a third
stage of substantive work, mandating one shared order at every
stage can cost up to $12.5\%$ of the synchronized makespan
(equivalently, the synchronized completion can run about
$14\%$ over the optimum). Whether deeper counterexample chains
push the constant beyond $1/8$ is open
(Section~\ref{sec:discussion}).
\end{remark}

%% ============================================================
\section{Computational falsification of the diamond/FJFJ
conjecture}
\label{sec:falsification}

The extended preprint of \citet{bowmanA} conjectured makespan
dominance for diamond
($S \to \{A, B\} \to T$) and fork-join--fork-join
($\{A,B\} \to C \to \{D,E\} \to F$) topologies, supported by
more than $2 \times 10^7$ verified instances. The conjecture
is false; Corollary~\ref{cor:characterization} replaces it
(both topologies contain three-stage paths). The earlier
evidence had two blind spots: the exhaustive $n = 2$ sweeps
varied works only at non-source stages (source works held
fixed), and the random $n = 3$ sweep diluted over a
$12$-dimensional work cube. Embedding the chain counterexample
of Proposition~\ref{prop:chain-failure} with adversarial
\emph{source} works produces failures immediately
(Table~\ref{tab:falsification}; all rows verified by
exhaustive enumeration of all per-stage priority orders, with
$k_v = 2$ throughout).

\begin{table}[h]
\centering
\caption{Embedded-chain counterexamples in topologies
previously conjectured safe. Works are listed per stage as
(job~1, job~2[, job~3]).}
\label{tab:falsification}
\begin{tabular}{llcc}
\toprule
\textbf{Topology} & \textbf{Works (stage: per job)} &
\textbf{Sync.\ best} & \textbf{Optimum} \\
\midrule
Diamond, $n{=}2$ & $A(1,2) \mid B(4,1) \mid C(1,1) \mid D(1,2)$
  & $4.0$ & $3.5$ \\
Diamond, $n{=}3$ & as above; job 3 has work 1 throughout
  & $4.5$ & $4.0$ \\
Wide diamond & chain on one branch; two weak branches
  & $4.0$ & $3.5$ \\
FJFJ & chain works on $A \to C \to D$; $B$, $E$, $F$ weak
  & $4.5$ & $4.0$ \\
\bottomrule
\end{tabular}
\end{table}

Beyond the constructed instances, the failure density is
substantial once source works vary. An exhaustive $n = 2$
diamond sweep over \emph{all} works in $\{1, \ldots, 4\}^8$
(including sources; $65{,}536$ instances) finds $362$
makespan-dominance failures ($0.55\%$). At $n = 3$ with works
sampled uniformly from $\{1, \ldots, 8\}$, random sampling
finds failures at a ${\sim}5\%$ rate ($2{,}497$ of
$50{,}000$). Moreover, the preprint's $n \geq 3$ diamond claim
is not reproducible: rerunning the archived sweep generator
with its original seed produces dominance failures within
$200{,}000$ instances. We therefore withdraw that empirical
claim along with the conjecture it supported; a correction
notice accompanies the archived preprint. The episode is
itself instructive: computational evidence gathered over a
high-dimensional parameter space can concentrate, unnoticed,
on a lower-dimensional slice of it.

Conversely, the positive side of the characterization is
corroborated beyond single sinks: an exhaustive sweep of a
depth-2 two-sink topology ($\{A, B\}$ sources;
$T_1 \leftarrow A$; $T_2 \leftarrow A, B$; all works in
$\{1, \ldots, 4\}^8$) finds \emph{zero} makespan-dominance
failures in $65{,}536$ instances, as
Theorem~\ref{thm:positive} requires---while \emph{pointwise}
domination fails in $35.7\%$ of them (sink by sink: some
schedule and sink admit no synchronized schedule dominating
every job's completion at that sink), as
Remark~\ref{rem:pointwise} leads one to expect.

%% ============================================================
\section{The flowtime boundary: concordant dominance for two
jobs}
\label{sec:flowtime}

The makespan boundary of
Corollary~\ref{cor:characterization} leaves open what
synchronization costs for \emph{flowtime} objectives on deep
chains. Proposition~\ref{prop:chain-failure} shows weighted
flowtime dominance fails at $L = 3$ in general. This section
proves the positive complement: for two jobs whose work
profiles are \emph{concordant}, the synchronized
lighter-job-first schedule minimizes total flowtime on chains
of every length, and this is sharp---it fails for three
concordant jobs and fails without concordance.

Throughout, an $L$-stage chain has stages $1, \ldots, L$ with
$k_\ell \geq 1$ workers; two jobs $A$ and $B$ have works
$w_\ell(A), w_\ell(B)$ at stage $\ell$, with processing times
$p_\ell = w_\ell(A)/k_\ell$ and $q_\ell = w_\ell(B)/k_\ell$.
Profiles are \emph{concordant} if $p_\ell \leq q_\ell$ for
every $\ell$ ($A$ is always the lighter job). A
\emph{per-stage schedule} is a vector
$\sigma \in \{A, B\}^L$ giving the priority job at each stage;
the \emph{synchronized} schedule is
$\sigma^* = (A, \ldots, A)$. Jobs are released to stage~1 at
time~0, and $C_\ell^\sigma(\cdot)$ denotes stage-$\ell$
completion times.

\begin{theorem}[Concordant flowtime dominance, $n = 2$]
\label{thm:concordant}
For any $L$-stage concordant chain with two jobs, the
synchronized schedule $\sigma^*$ minimizes total flowtime:
$C_L^{\sigma^*}(A) + C_L^{\sigma^*}(B) \leq
 C_L^{\sigma}(A) + C_L^{\sigma}(B)$
for every per-stage schedule $\sigma$.
\end{theorem}

The proof occupies the rest of this section. We may assume
$p_\ell > 0$; if $p_\ell = 0$ at some stage, job $A$ passes
through it instantly under every schedule and the stage can be
deleted.

\subsection{Single-stage transitions}

At a single stage with processing times $p \leq q$ and arrival
times $a_A, a_B$, only the first-arriving job is available
between the two arrivals and is processed regardless of
priority. The completion times in the cases we need are as
follows; in each overlap case the server is busy continuously
from the first arrival to the last completion.

\begin{lemma}[Single-stage transitions]
\label{lem:transitions}
At a single stage with $p \leq q$:

\smallskip\noindent
\textbf{Case 1:} $a_A \leq a_B$.
\emph{(1a) No overlap} ($a_A + p \leq a_B$): both priorities
give $C(A) = a_A + p$, $C(B) = a_B + q$.
\emph{(1b) Overlap, $A$-first} ($a_A + p > a_B$):
$C(A) = a_A + p$, $C(B) = a_A + p + q$.
\emph{(1c) Overlap, $B$-first} ($a_A + p > a_B$):
$C(B) = a_B + q$, $C(A) = a_A + p + q$.

\smallskip\noindent
\textbf{Case 2:} $a_B < a_A$.
\emph{(2a) No overlap} ($a_B + q \leq a_A$): both priorities
give $C(B) = a_B + q$, $C(A) = a_A + p$.
\emph{(2b) Overlap, $A$-first} ($a_B + q > a_A$):
$C(A) = a_A + p$, $C(B) = a_B + p + q$.
\emph{(2c) Overlap, $B$-first} ($a_B + q > a_A$):
$C(B) = a_B + q$, $C(A) = a_B + p + q$.
\end{lemma}

\begin{lemma}[Single-stage flowtime optimality]
\label{lem:single-stage}
At a single stage with $p \leq q$ and $a_A \leq a_B$:
\textup{(i)} $A$-first minimizes $C(A) + C(B)$;
\textup{(ii)} under $A$-first, $C(A) \leq C(B)$.
\end{lemma}

\begin{proof}
No overlap: both priorities coincide and
$C(A) = a_A + p \leq a_B + q = C(B)$. Overlap: the $A$-first
sum is $2a_A + 2p + q$ and the $B$-first sum is
$a_A + a_B + p + 2q$; the difference
$(a_A - a_B) + (p - q) \leq 0$. Under $A$-first,
$C(A) = a_A + p < a_A + p + q = C(B)$.
\end{proof}

\subsection{The sync value function}

Under $\sigma^*$ with initial arrivals $(\alpha, \beta)$,
$\alpha \leq \beta$, job $A$ is never preempted and completes
each stage first (Lemma~\ref{lem:single-stage}(ii)
inductively), giving
$C_\ell^*(A) = \alpha + \sum_{s \leq \ell} p_s$ and
$C_\ell^*(B) = \max(C_\ell^*(A), C_{\ell-1}^*(B)) + q_\ell$
with $C_0^*(B) = \beta$.

\begin{definition}[Sync value]
\label{def:Phi}
For an $L$-stage concordant chain and arrivals
$(\alpha, \beta)$ with $\alpha \leq \beta$, let
$\Phi_L(\alpha, \beta) = C_L^*(A) + C_L^*(B)$, so that
\[
  \Phi_L(\alpha, \beta)
  = \Phi_{L-1}\bigl(\alpha + p_1,\;
    \max(\alpha + p_1, \beta) + q_1\bigr),
  \qquad
  \Phi_0(\alpha, \beta) = \alpha + \beta ,
\]
where $\Phi_{L-1}$ refers to the chain with stage~1 removed.
\end{definition}

\begin{lemma}[Monotonicity]
\label{lem:phi-mono}
$\Phi_L(\alpha, \beta)$ is non-decreasing in both arguments
\textup{(}for $\alpha \leq \beta$\textup{)}.
\end{lemma}

\begin{proof}
Induction on $L$. For $L = 0$, $\Phi_0 = \alpha + \beta$. For
$L \geq 1$, $\Phi_L = \Phi_{L-1}(f, g)$ with
$f = \alpha + p_1$ and
$g = \max(\alpha + p_1, \beta) + q_1$ both non-decreasing in
$(\alpha, \beta)$, and $f \leq g$; conclude by the inductive
hypothesis.
\end{proof}

\begin{lemma}[Sum-preserving monotonicity]
\label{lem:sum-preserving}
For fixed $T = \alpha + \beta$ with $\alpha \leq \beta$,
$\Phi_L(\alpha, T - \alpha)$ is non-decreasing in $\alpha$.
\end{lemma}

\begin{proof}
Induction on $L$. For $L = 0$ the value is the constant $T$.
For $L \geq 1$,
$\Phi_L(\alpha, T - \alpha) = \Phi_{L-1}(f(\alpha),
g(\alpha))$ with $f(\alpha) = \alpha + p_1$ increasing and
$g(\alpha) = \max(\alpha + p_1, T - \alpha) + q_1$ piecewise
linear: $g$ decreases while $2\alpha + p_1 \leq T$ and
increases thereafter. On the decreasing regime,
$f + g = T + p_1 + q_1$ is constant with $f$ increasing, so
the inductive hypothesis (applied at sum $T + p_1 + q_1$)
gives the claim; on the increasing regime both arguments are
non-decreasing and Lemma~\ref{lem:phi-mono} applies.
\end{proof}

\subsection{Two interlocking inductions}

\begin{proposition}[$A$-first arrivals]
\label{prop:Pmain}
For any $L$-stage concordant chain, arrivals
$(\alpha, \beta)$ with $\alpha \leq \beta$, and any schedule
$\sigma$:
$\Phi_L(\alpha, \beta) \leq C_L^\sigma(A) + C_L^\sigma(B)$.
\end{proposition}

\begin{lemma}[$B$-first arrivals]
\label{lem:Qsecondary}
For any $L$-stage concordant chain, arrivals $(a, b)$ with
$a > b \geq 0$ \textup{(}$A$ arrives at $a$, $B$ at
$b$\textup{)}, and any schedule $\sigma$:
$C_L^\sigma(A) + C_L^\sigma(B) \geq \Phi_L(b, a)$.
\end{lemma}

Theorem~\ref{thm:concordant} is
Proposition~\ref{prop:Pmain} with $\alpha = \beta = 0$.
Lemma~\ref{lem:Qsecondary} says the ``wrong'' arrival order
cannot beat the sync value computed from the sorted arrivals.
We prove both simultaneously by induction on $L$.

\begin{proof}[Proof of Proposition~\ref{prop:Pmain} and
Lemma~\ref{lem:Qsecondary}]
\textbf{Base case $L = 0$.} Both sides equal
$\alpha + \beta$ (resp.\ $a + b$).

\textbf{Base case $L = 1$.} For
Proposition~\ref{prop:Pmain}, Lemma~\ref{lem:single-stage}(i)
states exactly that $A$-first minimizes the single-stage sum.
For Lemma~\ref{lem:Qsecondary} ($a > b$), check the cases of
Lemma~\ref{lem:transitions}:
\emph{(2a)} the sum is $a + b + p_1 + q_1$; since
$b + p_1 \leq b + q_1 \leq a$,
$\Phi_1(b, a) = (b + p_1) + (a + q_1)$, with equality.
\emph{(2b)} the sum is $a + b + 2p_1 + q_1$. If
$a \geq b + p_1$ then
$\Phi_1(b,a) = a + b + p_1 + q_1 \leq$ the sum; if
$a < b + p_1$ then
$\Phi_1(b,a) = 2b + 2p_1 + q_1 \leq$ the sum since $b < a$.
\emph{(2c)} the sum is $2b + p_1 + 2q_1$. If
$a \geq b + p_1$ then
$\Phi_1(b,a) = a + b + p_1 + q_1 \leq 2b + p_1 + 2q_1$ since
$a - b \leq q_1$ (overlap); if $a < b + p_1$ then
$\Phi_1(b,a) = 2b + 2p_1 + q_1 \leq 2b + p_1 + 2q_1$ since
$p_1 \leq q_1$.

\textbf{Inductive step for Proposition~\ref{prop:Pmain}.}
Arrivals $(\alpha, \beta)$, $\alpha \leq \beta$; stage-1
priority $\pi_1$. Recall
$\Phi_L(\alpha,\beta) = \Phi_{L-1}(\alpha + p_1,
\max(\alpha + p_1, \beta) + q_1)$.

\emph{Case $\pi_1 = A$.} Stage-1 outputs are
$C_1(A) = \alpha + p_1 \leq C_1(B) =
\max(\alpha + p_1, \beta) + q_1$ (cases 1a/1b), and the
$(L{-}1)$-stage Proposition~\ref{prop:Pmain} gives
$C_L^\sigma \geq \Phi_{L-1}(C_1(A), C_1(B))
= \Phi_L(\alpha, \beta)$.

\emph{Case $\pi_1 = B$.} Without overlap
($\alpha + p_1 \leq \beta$) the outputs coincide with the
$A$-first outputs and the previous case applies. With overlap,
outputs are $C_1(A) = \alpha + p_1 + q_1 >
C_1(B) = \beta + q_1$ (case 1c), so the $(L{-}1)$-stage
Lemma~\ref{lem:Qsecondary} gives
$C_L^\sigma \geq \Phi_{L-1}(\beta + q_1,\,
\alpha + p_1 + q_1)$. Since overlap means
$\alpha + p_1 > \beta$,
$\Phi_L(\alpha, \beta) = \Phi_{L-1}(\alpha + p_1,\,
\alpha + p_1 + q_1)$; the second arguments agree, and
$\alpha + p_1 \leq \beta + q_1$ (from $\alpha \leq \beta$,
$p_1 \leq q_1$), so Lemma~\ref{lem:phi-mono} in the first
argument completes the case.

\textbf{Inductive step for Lemma~\ref{lem:Qsecondary}.}
Arrivals $(a, b)$, $a > b$; stage-1 priority $\pi_1$. Write
$\mu = \max(b + p_1, a)$, so
$\Phi_L(b, a) = \Phi_{L-1}(b + p_1, \mu + q_1)$.

\emph{Case 2a (no overlap, $b + q_1 \leq a$).} Both
priorities give $C_1(A) = a + p_1 > C_1(B) = b + q_1$, and the
$(L{-}1)$-stage Lemma~\ref{lem:Qsecondary} gives
$C_L^\sigma \geq \Phi_{L-1}(b + q_1, a + p_1)$. Here
$\mu = a$ (as $b + p_1 \leq b + q_1 \leq a$), so
$\Phi_L(b,a) = \Phi_{L-1}(b + p_1, a + q_1)$. The pairs
$(b + q_1, a + p_1)$ and $(b + p_1, a + q_1)$ share the sum
$a + b + p_1 + q_1$, with
$b + p_1 \leq b + q_1 \leq a + p_1$, so
Lemma~\ref{lem:sum-preserving} gives
$\Phi_{L-1}(b + q_1, a + p_1) \geq
 \Phi_{L-1}(b + p_1, a + q_1)$.

\emph{Case 2b (overlap, $\pi_1 = A$).} Outputs
$C_1(A) = a + p_1 < C_1(B) = b + p_1 + q_1$ (using
$b + q_1 > a$), so the $(L{-}1)$-stage
Proposition~\ref{prop:Pmain} gives
$C_L^\sigma \geq \Phi_{L-1}(a + p_1,\, b + p_1 + q_1)$.
If $a \geq b + p_1$ ($\mu = a$): the pairs
$(a + p_1, b + p_1 + q_1)$ and $(b + p_1, a + p_1 + q_1)$
share the sum $a + b + 2p_1 + q_1$, each has first argument
$\leq$ second (for the first pair because $a < b + q_1$), and
$a + p_1 \geq b + p_1$; Lemma~\ref{lem:sum-preserving} then
Lemma~\ref{lem:phi-mono} (second argument,
$a + p_1 + q_1 \geq a + q_1$) give
\[
  \Phi_{L-1}(a {+} p_1, b {+} p_1 {+} q_1)
  \geq \Phi_{L-1}(b {+} p_1, a {+} p_1 {+} q_1)
  \geq \Phi_{L-1}(b {+} p_1, a {+} q_1) = \Phi_L(b, a).
\]
If $a < b + p_1$ ($\mu = b + p_1$): by
Lemma~\ref{lem:phi-mono} in the first argument
($a + p_1 > b + p_1$),
$\Phi_{L-1}(a + p_1, b + p_1 + q_1) \geq
 \Phi_{L-1}(b + p_1, b + p_1 + q_1) = \Phi_L(b, a)$.

\emph{Case 2c (overlap, $\pi_1 = B$).} Outputs
$C_1(B) = b + q_1 < C_1(A) = b + p_1 + q_1$, so the
$(L{-}1)$-stage Lemma~\ref{lem:Qsecondary} gives
$C_L^\sigma \geq \Phi_{L-1}(b + q_1,\, b + p_1 + q_1)$.
If $a < b + p_1$ ($\mu = b + p_1$): monotonicity in the first
argument ($b + q_1 \geq b + p_1$) gives
$\Phi_{L-1}(b + q_1, b + p_1 + q_1) \geq
 \Phi_{L-1}(b + p_1, b + p_1 + q_1) = \Phi_L(b, a)$.
If $a \geq b + p_1$ ($\mu = a$), split on
$b + q_1 \lessgtr a + p_1$. If $b + q_1 \leq a + p_1$:
monotonicity in the second argument
($b + p_1 + q_1 \geq a + p_1$, from $b + q_1 > a$) gives
$\Phi_{L-1}(b + q_1, b + p_1 + q_1) \geq
 \Phi_{L-1}(b + q_1, a + p_1)$, and the pairs
$(b + q_1, a + p_1)$, $(b + p_1, a + q_1)$ share their sum
with $b + p_1 \leq b + q_1 \leq a + p_1$, so
Lemma~\ref{lem:sum-preserving} finishes. If
$b + q_1 > a + p_1$: by Lemma~\ref{lem:phi-mono}
(componentwise, since $b + q_1 \geq a + p_1$ and
$b + p_1 + q_1 \geq b + q_1$),
$\Phi_{L-1}(b + q_1, b + p_1 + q_1) \geq
 \Phi_{L-1}(a + p_1, b + q_1)$; the pairs
$(a + p_1, b + q_1)$ and $(b + p_1, a + q_1)$ share their sum
with $a + p_1 \geq b + p_1$ and first $\leq$ second in each,
so Lemma~\ref{lem:sum-preserving} gives
$\Phi_{L-1}(a + p_1, b + q_1) \geq
 \Phi_{L-1}(b + p_1, a + q_1) = \Phi_L(b, a)$.

This closes both inductions and proves
Theorem~\ref{thm:concordant}.
\end{proof}

\begin{remark}[Sharpness]
\label{rem:flowtime-sharp}
Concordance is essential, already for unweighted flowtime: for
the discordant works $w = [(1,6),(6,1),(6,1)]$ with $k = 2$,
the best synchronized total flowtime is $13.5$ while the
per-stage orders $(1,2) \mid (2,1) \mid (2,1)$ attain $11.5$
(and Proposition~\ref{prop:chain-failure}'s instance fails for
weighted flowtime). Two jobs are essential as well: for $L = 2$,
$n = 3$, $k = 2$, $w = [(4,5,6),\,(1,9,10)]$ (strictly
concordant), the synchronized lighter-first schedule attains
total flowtime $25.5$, while the per-stage orders
$(2,1,3) \mid (1,2,3)$ attain $25.0$ by letting the lightest
job preempt a medium job at stage~2. In an exhaustive sweep
this failure is rare: it is the sole counterexample among
$14{,}400$ strictly concordant $n = 3$, $L = 2$ configurations
with work values $1$--$10$, and $175{,}616$ strictly
concordant $n = 3$, $L = 3$ configurations with work values
$1$--$8$ contain no flowtime failure. The $n = 2$ theorem is
itself verified exhaustively for $L \in \{2, 3\}$, $k = 2$,
work values $1$--$8$: all $784$ strictly concordant two-stage
and $21{,}952$ strictly concordant three-stage configurations
($46{,}656$ at $L = 3$ when ties are allowed) show zero
flowtime failures.
\end{remark}

%% ============================================================
\section{Practice: policy benchmarks}
\label{sec:practice}

The negative results say that at depth three or more, full
synchronization is no longer free. They do not say what a
practitioner should do. This section evaluates the policy we
recommend and quantifies both its aggregate cost and the
structure of the residual gap. All statistics below are
produced by exhaustive or seeded-random benchmarks in the
companion artifact (see Code availability); the model is that
of Definition~\ref{def:model}.

\begin{definition}[Composite policy: source-sync with FIFO
downstream]
\label{def:composite}
Fix one order $\pi$. Every source stage uses priority $\pi$.
Every non-source stage serves first-in-first-out by arrival
time, breaking ties by $\pi$.
\end{definition}

The policy asks for global agreement only where work
originates; downstream teams need no ranking at all---they
simply take work in the order it reaches them. On
single-source topologies this is provably not a relaxation of
full synchronization but a reimplementation of it:

\begin{lemma}[Propagation of source order]
\label{lem:propagation}
Let $G$ have a single source. Then for every work assignment,
the composite policy with order $\pi$ produces exactly the
completion times of the fully synchronized schedule
$\beta_v \equiv \pi$, at every stage.
\end{lemma}

\begin{proof}
By induction over a topological order, every stage's arrival
vector is weakly $\pi$-ordered and its effective service order
is $\pi$. At the source, all jobs arrive at time~0 and the
policy uses $\pi$ by definition; by
Lemma~\ref{lem:prefix}(v), completions are $\pi$-ordered. At a
non-source stage $v$, each predecessor's completions are
$\pi$-ordered by induction, so the componentwise maximum
defining $a_v$ is weakly $\pi$-ordered; FIFO-by-arrival with
ties broken by $\pi$ therefore sorts jobs exactly in the order
$\pi$. With $\pi$-ordered arrivals and priority $\pi$, no job
ever overtakes another in service
(Lemma~\ref{lem:prefix}(v)): completions follow the max-plus
recurrence and are again $\pi$-ordered---and they coincide
with the synchronized schedule's completions at $v$, since
that schedule has the same arrivals (inductively) and the same
priority order.
\end{proof}

\begin{remark}
\label{rem:multisource}
The argument extends verbatim to multi-source DAGs when all
sources share the order $\pi$ and all jobs are released
simultaneously, since componentwise maxima of $\pi$-ordered
vectors are $\pi$-ordered. The benchmarks below corroborate
this exactly: per order $\pi$, the composite policy and full
synchronization produced identical sink makespans on all
$561{,}633$ chain and fork-join--fork-join instances tested
(per-$\pi$ mismatch count zero)---the chain and
fork-join--fork-join families of Table~\ref{tab:policy} plus
a $4{,}096$-instance exhaustive fork-join--fork-join sweep
with works $1$--$2$---including, in particular, the
two-source fork-join--fork-join topology. Downstream
discipline is a non-issue once the sources are aligned.
\end{remark}

\paragraph{Aggregate quality.}
On the exhaustive diamond family of
Section~\ref{sec:falsification} ($n = 2$, $k_v = 2$, all works
in $\{1,\ldots,4\}^8$; $65{,}536$ instances), with the source
order chosen by enumeration and compared against the
enumerated optimum over all per-stage orders, the composite
policy is optimal for the makespan in $99.45\%$ of instances,
with mean gap $0.058\%$ and maximum gap $14.29\%$; for total
flowtime at the sink it is optimal in $99.41\%$ of instances,
with mean gap $0.036\%$ and maximum gap $11.76\%$.

\paragraph{Across topology families.}
The same pattern holds beyond the diamond.
Table~\ref{tab:policy} reports the identical benchmark on
exhaustive and seeded-random families over three-stage chains
and the two-source fork-join--fork-join topology of
Section~\ref{sec:falsification} ($k_v = 2$ throughout). In
every family the policy's optimality rate and full
synchronization's failure rate sum to one: by the per-$\pi$
identity of Remark~\ref{rem:multisource} the composite policy
is optimal exactly where full synchronization is, so it
inherits full synchronization's failures and nothing else.
Failure rates grow with the work range and with $n$---from
$0.55\%$ on the exhaustive diamond family to $10.05\%$ on
random three-job fork-join--fork-join instances---but the
conditional cost where synchronization fails stays in one
band: mean gap $5.5$--$10.4\%$, maximum $12.5$--$20\%$.

\begin{table}[h]
\centering
\caption{Composite policy against the enumerated optimum
across topology families (chains are three-stage; ranges are
integer work values; exh.\ $=$ exhaustive, rnd.\ $=$ seeded
random). ``Pol.\ opt.''\ is the fraction of instances where
the best-$\pi$ policy makespan equals the optimum; ``sync
fails'' is the fraction where the best fully synchronized
makespan is suboptimal (the two columns sum to $100\%$); the
gap columns describe the relative excess over the optimum on
the failing instances.}
\label{tab:policy}
\small
\setlength{\tabcolsep}{4pt}
\begin{tabular}{lrrrrr}
\toprule
\textbf{Family} & \textbf{Inst.} &
\textbf{Pol.\ opt.} & \textbf{Sync fails} &
\textbf{Gap mean} & \textbf{Gap max} \\
\midrule
Diamond $n{=}2$, $1$--$4$ (exh.) & $65{,}536$
  & $99.45\%$ & $0.55\%$ & $10.5\%$ & $14.3\%$ \\
Chain $n{=}2$, $1$--$4$ (exh.) & $4{,}096$
  & $99.07\%$ & $0.93\%$ & $10.4\%$ & $14.3\%$ \\
Chain $n{=}3$, $1$--$6$ (rnd.) & $20{,}000$
  & $94.03\%$ & $5.97\%$ & $6.3\%$ & $20.0\%$ \\
FJFJ $n{=}2$, $1$--$3$ (exh.) & $531{,}441$
  & $98.68\%$ & $1.32\%$ & $8.9\%$ & $12.5\%$ \\
FJFJ $n{=}3$, $1$--$6$ (rnd.) & $2{,}000$
  & $89.95\%$ & $10.05\%$ & $5.5\%$ & $15.0\%$ \\
\bottomrule
\end{tabular}
\end{table}

\paragraph{Cost of synchronization where it fails.}
On the same family, full synchronization is makespan-suboptimal
on $362$ instances ($0.55\%$); on those instances the gap to
the optimum has mean $10.5\%$, median $10.0\%$, and maximum
$14.3\%$. This converts the qualitative ``dominance can
fail'' of Theorem~\ref{thm:negative} into an expected cost:
failures are rare but material when present. The maximum is
attained by the embedded chain instance of
Table~\ref{tab:falsification} (synchronized $4$ against
optimal $3.5$: $1/7 \approx 14.3\%$ above the optimum,
i.e.\ the optimum is $1/8$ below the synchronized makespan,
matching the asymptotic constant of
Theorem~\ref{thm:negative}). By
Lemma~\ref{lem:propagation}, the composite policy coincides
with full synchronization on this single-source family, so its
worst instances are the same $0.55\%$: capturing the residual
$10$--$14\%$ there requires deliberate priority
overtaking---expediting the work-inverted job---which no
arrival-order discipline reproduces.

\paragraph{Robustness to estimation error.}
Rankings are made on estimates. On diamonds with $n = 3$
($5{,}000$ seeded instances; estimated works drawn from
$\{10, 20, \ldots, 60\}$; true works equal to the estimate
scaled by an independent uniform factor on $[0.5, 1.5]$, i.e.\
up to $\pm 50\%$ error), choosing $\pi$ from the
\emph{estimates}, executing on the \emph{true} works with FIFO
adapting to actual arrivals at runtime, and measuring against
the clairvoyant optimum on true works: the composite policy is
optimal outright in $50.1\%$ of instances, with mean gap
$3.03\%$ and maximum $32.5\%$. The picture is stable across
topologies: under the same protocol the mean gap is $3.17\%$
on three-stage chains ($5{,}000$ instances; optimal in
$50.9\%$, maximum $33.6\%$) and $4.11\%$ on the two-source
fork-join--fork-join family ($1{,}000$ instances; optimal in
$38.8\%$, maximum $29.8\%$); in all three families the
executed policy and the executed fully synchronized schedule
again had identical makespans on every instance. The bulk of
alignment's benefit survives severe estimation error; the
residual $3$--$4\%$ is the price of uncertainty itself, not of
the policy.

\paragraph{Managerial reading.}
Agree on one ranked list where work enters the system; let
every downstream team serve work in the order it arrives.
On single-source delivery structures this reproduces the
fully synchronized schedule automatically---no further
coordination meetings, no per-team re-ranking---and work
reaches the final stage as soon as it would under
organization-wide synchronization. The residual upside over
that baseline is rare ($0.55$--$10\%$ of instances across the
families of Table~\ref{tab:policy}), worth roughly
$5$--$10\%$ when present, and lives in identifiable
cross-stage work inversions: an item that is light for the
upstream team but heavy downstream
(Remark~\ref{rem:mechanism}). The playbook is therefore:
align the source; let arrival order govern everything below
it; and treat inversions as named exceptions to expedite
deliberately, rather than re-sequencing every team's queue.

%% ============================================================
\section{Discussion}
\label{sec:discussion}

Table~\ref{tab:landscape} assembles the corrected dominance
landscape. The single organizing fact is
Corollary~\ref{cor:characterization}: for the makespan, the
boundary of safe synchronization is depth, not width. The
diamond and fork-join rows of our earlier preprint's landscape
table claimed computational support for dominance on those
topologies; both rows are corrected here, and the depth
criterion subsumes the would-be taxonomy of tree and
layered-DAG special cases (any such topology is safe precisely
when it contains no directed three-stage path).

\begin{table}[h]
\centering
\caption{Dominance landscape (corrected). All entries proved
in this paper or in \citet{bowmanA}.}
\label{tab:landscape}
\begin{tabular}{lll}
\toprule
\textbf{Model} & \textbf{Property} & \textbf{Status} \\
\midrule
Two-stage assembly \citep{bowmanA}
  & Pointwise (downstream inherited)
  & Holds \\
Depth-2 DAG, any sink (Thm.~\ref{thm:positive})
  & Makespan dominance
  & Holds \\
Chain, $L = 2$ (Thm.~\ref{thm:chain-l2})
  & Makespan dominance
  & Holds \\
Chain, $L = 2$ (Rem.~\ref{rem:pointwise})
  & Pointwise (fully synchronized)
  & Fails \\
Chain, $L \geq 3$ (Prop.~\ref{prop:chain-failure})
  & Makespan / weighted flowtime
  & Fails \\
DAG with a 3-path (Thm.~\ref{thm:negative})
  & Makespan dominance
  & Fails \\
Concordant chain, $n = 2$ (Thm.~\ref{thm:concordant})
  & Total flowtime
  & Holds \\
Concordant chain, $n \geq 3$ (Rem.~\ref{rem:flowtime-sharp})
  & Total flowtime
  & Fails \\
\bottomrule
\end{tabular}
\end{table}

One adjacent question is settled computationally: the
depth-three failures \emph{survive} the integer,
non-preemptive divisible-work model of \citet{bowmanA}, but
only barely. Chaining that model through three stages ($k = 2$
workers per stage; each job's integer stage work split freely
into integer per-worker portions; no preemption), an
exhaustive sweep of all $4{,}096$ two-job instances with stage
works at most $4$ finds exactly two makespan-dominance
failures---a single instance up to job relabeling,
$w = [(1,3),\,(4,1),\,(1,3)]$, with synchronized best $6$
against optimum $5$; the optimum requires both an order
inversion at the final stage and an uneven integer split of
the heavy middle-stage work. The failure is a granularity
knife-edge: scaling that instance's works by $2$, $3$, or $4$
closes the gap entirely (synchronized $=$ optimum
$= 9, 15, 18$), and the fluid counterexample of
Proposition~\ref{prop:chain-failure} never fails in the
integer model at any tested scale (work multipliers up
to~$8$), even though its fluid gap grows linearly.
Indivisibility removes the mid-stage overtaking the fluid
counterexamples exploit, so the depth boundary is real but
attenuated by integrality: it bites only at work
granularities comparable to the worker count.

Two questions remain open. First, whether deeper embedded
chains can push the asymptotic optimality gap beyond $1/8$ of
the synchronized makespan. Second, stochastic processing
times: the robustness benchmark of
Section~\ref{sec:practice} suggests the source-alignment
policy degrades gracefully under estimation error, but a
distributional theory of synchronization dominance is
untouched.

%% ============================================================

\section*{Declarations}

\textbf{Funding.}
This research did not receive any specific grant from funding
agencies in the public, commercial, or not-for-profit sectors.

\textbf{Competing interests.}
The author declares no competing interests.

\textbf{Author contributions.}
The sole author conceived the work, developed the proofs, and
wrote the manuscript.

\textbf{Data availability.}
The verification and benchmark data supporting the
quantitative claims of
Sections~\ref{sec:falsification}--\ref{sec:practice} are
available in the repository below.

\textbf{Code availability.}
The simulation and benchmark code is available at
\url{https://github.com/ebowman/synchronization-paper-artifact}
and archived at
\url{https://doi.org/10.5281/zenodo.20075388}.

\textbf{Generative AI disclosure.}
During the preparation of this work, the author used Claude
for critique of exposition, structural feedback, and language
revision. The author reviewed and edited the manuscript and
takes full responsibility for the content.

%% ============================================================

\bibliography{boundary}

\end{document}
